\writeSectionFrame{\presentationTitle}{Dokumentenbeispiel}
\begin{frame}{Klassendiagramm Dokumente ohne Brücke}
	\begin{center}
 		\includegraphics[scale=0.6]{bilder/dokumente1.png}
 	\end{center}
\end{frame}

\note{
	Wir sind noch ein etwas realistischeres Beispiel an. Unsere Anwendung soll
	verschiedene Dokumente verarbeiten. Diese Dokumente können unterschiedlich
	gespeichert werden. Wir können eine FileStream API zum Lesen und Schreiben
	verwenden. Aber es gibt Tochterunternehmen, die noch mit allen Servern unterwegs
	sind, die diese API nicht unterstützen. Und dann gibt es noch Archivsysteme, wo
	die Dokumente u.U. auch untergebracht werden müssen. Wir haben also vier Dokumentenklassen
	und drei Speichersysteme und bräuchten damit 12 Klassen in einer Vererbungshierarchie. 
	Und das nur, wenn es bei diesen vier Dokumentenklassen bleibt.  
}

\begin{frame}{Klassendiagramm Dokumente mit Brücke}
	\begin{center}
 		\includegraphics[scale=0.35]{bilder/dokumente2.png}
 	\end{center}
\end{frame}

\begin{frame}{Abstraktion}
	\begin{exampleblock}{Document.java}
		\includegraphics[scale=0.3]{bilder/dokumente3.png}
	\end{exampleblock}
\end{frame}

\begin{frame}{KonkreteAbstraktion}
	\begin{exampleblock}{Invoice.java}
		\includegraphics[scale=0.25]{bilder/dokumente4.png}
	\end{exampleblock}
\end{frame}

\begin{frame}{Implementierer}
	\begin{exampleblock}{PersistenceProvider.java}
		\includegraphics[scale=0.45]{bilder/dokumente5.png}
	\end{exampleblock}
\end{frame}

\begin{frame}{Implementierer}
	\begin{exampleblock}{FileSystemProvider.java}
		\includegraphics[scale=0.35]{bilder/dokumente6.png}
	\end{exampleblock}
\end{frame}